\documentclass[10pt,a4paper]{amsart}
\usepackage[margin=19mm]{geometry}
\usepackage{amsmath,amssymb,mathtools}
\usepackage[scr=rsfs,cal=euler]{mathalfa}
\usepackage{xfrac}
\usepackage[unicode,hidelinks,bookmarks=false]{hyperref}
\newcommand{\ie}{i.e.\ }
\newcommand{\cf}{cf.\ }
\newcommand{\resp}{respectively\ }
\newcommand{\Subsection}{\subsection}
\providecommand{\nobreakdash}{\nobreak\hbox{-}}
\setlength{\emergencystretch}{2em}
\multlinegap0pt
\usepackage{mathtools}
\usepackage{amssymb}
\usepackage{dsfont}
\usepackage{tikz}
\newtheorem{proposition}{Proposition}
\newtheorem{lemma}{Lemma}
\newtheorem{theorem}{Theorem}
\newcommand{\EC}{\mathcal{E}\mathcal{C}}
\newcommand{\edge}{\mathbf{e}}
\title{A decomposition of graph $a$-numbers}
\author{Suyuong Choi, Younghan Yoon}
\date{Source: arXiv:2508.06855v2 (CC BY 4.0)}
\begin{document}
\maketitle

\noindent\emph{Source and licence.} A decomposition of graph $a$-numbers. Version arXiv:2508.06855v2 (CC BY 4.0), distributed by arXiv under the Creative Commons Attribution 4.0 International licence, http://creativecommons.org/licenses/by/4.0/, as recorded in the arXiv licence metadata for this exact version and shown on the abstract page https://arxiv.org/abs/2508.06855v2. Full licence notice: \texttt{licenses/arxiv-2508.06855-CC-BY-4.0.txt}.\par\smallskip

\noindent\textit{Editorial scope.} Counted unit: the article's theorem thm:decom (source0.tex lines 903--909). The article's complete original proof of it is delivered with it (source0.tex lines 910--945, one \texttt{proof} environment of 36 lines). No mathematics is written, repaired or re-derived: the statement and the proof are the article's own lines, in the article's own notation, and no retained line is substituted or re-flowed.  Retained uncounted as the source-local context the proof needs: lines 382--384; lines 732--741; lines 854--860.  Nothing was omitted from any retained span, so the package carries no omission record.  No retained line contains a citation, so the wrapper carries no bibliography and no bibliography entry.  No retained line contains a figure environment, an image-inclusion command or drawing code, so no image is delivered and the package declares no asset dependency.  Editorial formatting only, in addition: an AMS \texttt{amsart} preamble that restates the article's own declarations and macros used by the retained lines, namely the environments \texttt{proposition}, \texttt{lemma}, \texttt{theorem} on separate counters, together with the standard packages those lines need.  The retained environments print in this document as Proposition 1, Lemma 1, Theorem 1 in order of appearance, whereas the article prints the same items with the numbering of its own full document.  The complete native source of the article as distributed in the arXiv e-print for this exact version is bundled unchanged as \texttt{sources/182970/source0.tex}.
\begin{equation}\label{eq:inter}
  (G\vert_{J})^\ast_{I} = G^\ast_{I}\vert_{J \setminus I}.
\end{equation}

\begin{proposition}\label{prop:conn}
    For a finite simple graph~$G$, let~$J \in \EC(G)$, and~$\edge$ a pair of vertices of~$G$.
    The following statements hold.
    \begin{enumerate}
      \item\label{prop:conn1} For~$I \subseteq J$, $I \in \EC(G)$ if and only if~$I \in \EC(G\vert_J)$. 
      \item\label{prop:conn2} For~$J \subseteq I$, $I \in \EC(G)$ if and only if~$I \setminus J \in \EC(G^\ast_J)$.
      \item\label{prop:conn3} $\EC(G) \subseteq \EC(G+\edge)$.
    Moreover, $I \in \EC(G+\edge) \setminus \EC(G)$ implies that $\edge \subseteq I$.
    \end{enumerate}
\end{proposition}

\begin{lemma}\label{lemma:decom}
    Let $G$ be a simple graph on~$2i$ vertices.
    For a pair~$\edge$ of vertices of~$G$,
    $$
    a(G+\edge) = a(G) + \sum_{J \in \EC(G+\edge) \setminus \EC(G)} \left\vert b(G\vert_J) \right\vert \cdot a((G+\edge)^\ast_J).
    $$
\end{lemma}

\begin{theorem}\label{thm:decom}
    Let $G$ be a finite simple graphs with the vertex set~$V(G)$, and $\edge$ a pair of vertices of~$G$.
For each $i \geq 0$, 
    \begin{equation} \label{eqn:main_decomp}
    a_i(G+\edge) - a_i(G) = \sum_{J \in \EC(G+\edge) \setminus \EC(G)}\left\vert  b(G\vert_J) \right\vert \cdot a_{i- \frac{|J|}{2}}((G +\edge)^\ast_{J}).
    \end{equation}
\end{theorem}

\begin{proof}
    To begin, since~$V(G) = V(G+\edge)$, we have
    $$a_i(G+\edge) = \sum_{\substack{I\subseteq V(G)\\ |I| =2i}}a((G+\edge)\vert_I).
    $$
    We now divide the sum according to whether~$I$ contains~$\edge$ or not; that is,
    $$
    a_i(G+\edge) = \sum_{\substack{\edge \not\subset I\subseteq V(G)\\ |I| =2i}} a((G+\edge)\vert_I) + \sum_{\substack{\edge \subseteq I\subseteq V(G)\\ |I| =2i}} a((G+\edge)\vert_I).
    $$
    For each~$I \in \EC(G + \edge)$, note that
    $$
    (G+\edge)\vert_I = \begin{cases}
      G\vert_I +\edge, & \mbox{if } \edge \in I \\
      G\vert_I, & \mbox{otherwise}.
    \end{cases}
    $$
    It follows that
    $$
    a_i(G+\edge) = \sum_{\substack{\edge \not\subset I\subseteq V(G)\\ |I| =2i}} a(G\vert_I) + \sum_{\substack{\edge \subseteq I\subseteq V(G)\\ |I| =2i}} a(G\vert_I+\edge).
    $$
    We then apply Lemma~\ref{lemma:decom} to~$G\vert_I +\edge$, which confirms that
    \begin{align*}
      a_i(G+\edge) & = \sum_{\substack{\edge \not\subset I\subseteq V(G)\\ |I| =2i}} a(G\vert_I) + \sum_{\substack{\edge \subseteq I\subseteq V(G)\\ |I| =2i}} \Big(a(G\vert_I) + \sum_{J \in \EC((G+\edge)\vert_I) \setminus \EC(G\vert_I)} \left\vert b(G\vert_J) \right\vert a\big(((G+\edge)\vert_I)^\ast_J\big)\Big) \\
       & = a_i(G) + \sum_{\substack{\edge \subseteq I\subseteq V(G)\\ |I| =2i}} \sum_{\substack{J \in \EC(G+\edge) \setminus \EC(G) \\ J \subseteq I}} \left\vert b(G\vert_J) \right\vert a\big(((G+\edge)\vert_I)^\ast_J\big) \quad \text{ (by Proposition~\ref{prop:conn}~\eqref{prop:conn1})}.
    \end{align*}
    From Proposition~\ref{prop:conn}~\eqref{prop:conn3}, for each~$J \in \EC(G+\edge) \setminus \EC(G)$, we have $\edge \subseteq J$.
    It follows that interchanging the order of summation in the second term on the right-hand side yields
    $$
       a_i(G+\edge) = a_i(G) + \sum_{J \in \EC(G+\edge) \setminus \EC(G)}\Big( \left\vert b(G\vert_J) \right\vert \sum_{\substack{J \subseteq I\\ |I| =2i}} a\big(((G+\edge)\vert_I)^\ast_J\big)\Big).
    $$
    By replacing $I' =I \setminus J$ with~\eqref{eq:inter}, we obtain that
    \begin{align*}
      a_i(G+\edge) & = a_i(G) + \sum_{J \in \EC(G+\edge) \setminus \EC(G)}\Big( \left\vert b(G\vert_J) \right\vert \sum_{\substack{I' \in V((G+\edge)^\ast_J)\\ |I'| =2i-\left\vert J \right\vert}} a\big(((G+\edge)^\ast_J)_{I'})\big)\Big) \\
       & = a_i(G) + \sum_{J \in \EC(G+\edge) \setminus \EC(G)} \left\vert b(G \vert_J) \right\vert \cdot a_{i-\frac{|J|}{2}}((G +\edge)^\ast_J),
    \end{align*}
as desired.
\end{proof}

\end{document}
